% Vignette of the Pollis method TwoScalesVolatility
% Compile with: julia -e 'using PollisTeX; display(compile("TwoScalesVolatility.tex"))'
% The look of a Pollis method vignette (the fonts and colours of the Pollis Article Methods documents).
% Set \vignetteauthor, \vignettedate and, if not CC BY-SA 4.0, \vignettelicense after % The look of a Pollis method vignette (the fonts and colours of the Pollis Article Methods documents).
% Set \vignetteauthor, \vignettedate and, if not CC BY-SA 4.0, \vignettelicense after % The look of a Pollis method vignette (the fonts and colours of the Pollis Article Methods documents).
% Set \vignetteauthor, \vignettedate and, if not CC BY-SA 4.0, \vignettelicense after \input{vignette_style}.
\documentclass[10pt,a4paper]{article}

\usepackage[a4paper,left=18mm,right=18mm,top=15mm,bottom=17mm,footskip=8mm]{geometry}
\usepackage{iftex}
\ifPDFTeX\usepackage[T1]{fontenc}\fi
\usepackage[default,semibold]{sourcesanspro}
\usepackage{textcomp}
\usepackage[italic]{mathastext}
\usepackage{amsmath}
\usepackage{microtype}
\usepackage[table]{xcolor}
\usepackage{tikz}
\usepackage[most]{tcolorbox}
\usepackage{fontawesome5}
\usepackage{tabularx,booktabs}
\usepackage{enumitem}
\usepackage{fancyhdr}
\usepackage[hidelinks]{hyperref}

% ---------------------------------------------------------------- colours (as in the brochure)
\definecolor{navy}{HTML}{1C3A5E}
\definecolor{pblue}{HTML}{3B78B5}
\definecolor{amber}{HTML}{F2B33D}
\definecolor{deepamber}{HTML}{D9902A}
\definecolor{pale}{HTML}{EEF3F9}
\definecolor{ink}{HTML}{2A2F36}
\definecolor{logonavy}{HTML}{173D63}
\definecolor{logodark}{HTML}{C98A1E}
\definecolor{logoamber}{HTML}{F2B446}
\definecolor{logocream}{HTML}{FEEFD2}
\definecolor{soft}{HTML}{5E6773}
\definecolor{agree}{HTML}{2F7D4F}
\definecolor{differ}{HTML}{B23B3B}
\color{ink}

\setlength{\parindent}{0pt}
\setlength{\parskip}{4pt}
\setlist{nosep,leftmargin=*,itemsep=2pt}
\renewcommand{\arraystretch}{1.15}

\pagestyle{fancy}
\fancyhf{}
\renewcommand{\headrulewidth}{0pt}
\fancyfoot[L]{\footnotesize\color{soft}\textbf{Pollis}\enspace Article Methods\enspace·\enspace\nmtheme}
\fancyfoot[R]{\footnotesize\color{soft}\thepage}
\newcommand{\nmtheme}{}
\newcommand{\vignetteauthor}{}
\newcommand{\vignettedate}{\today}
\newcommand{\vignettelicense}{\href{https://creativecommons.org/licenses/by-sa/4.0/}{CC BY-SA 4.0}}

% The Pollis logo, redrawn so it prints sharply at any size
\newcommand{\pollislogo}[1]{%
  \begin{tikzpicture}[x=#1/512,y=-#1/512,baseline=(current bounding box.center)]
    \fill[logonavy] (0,0) rectangle (512,512);
    \fill[logodark] (36,422) -- (104,256) -- (172,422) -- cycle;
    \fill[logoamber] (104,422) -- (203,180) -- (302,422) -- cycle;
    \fill[logocream] (203,422) -- (339,89) -- (476,422) -- cycle;
  \end{tikzpicture}}

% Top of the first page: logo, name; then the kicker, title and a one-line promise
\newcommand{\nmhead}[3]{%
  \renewcommand{\nmtheme}{#1}%
  \thispagestyle{fancy}%
  \begin{minipage}[c]{13mm}\pollislogo{11mm}\end{minipage}%
  \begin{minipage}[c]{\dimexpr\linewidth-13mm\relax}
    {\color{navy}\Large\bfseries Pollis}\enspace{\color{soft}Article Methods}\hfill
    {\color{deepamber}\small\bfseries\MakeUppercase{#1}}
  \end{minipage}\par\vspace{7mm}
  {\color{navy}\fontsize{22}{26}\selectfont\bfseries #2}\par\vspace{4pt}
  {\color{soft}\large #3}\par\vspace{5mm}}

\newcommand{\nmsection}[1]{\par\vspace{5pt}{\color{navy}\large\bfseries #1}\par\vspace{1pt}}
\newcommand{\nmsub}[1]{\par\vspace{2pt}{\color{deepamber}\bfseries #1}\par}
\newcommand{\gt}{\,{\color{deepamber}\bfseries ›}\,}
\newcommand{\menupath}[1]{{\bfseries\color{navy}#1}}
\newcommand{\fn}[1]{\texttt{\small #1}}
% Citations: \nmcite{key}{Breiman (2001)} links to the entry \nmref{key}{...} in the References
\newcommand{\nmcite}[2]{\hyperlink{ref:#1}{\color{pblue}#2}}
\newcommand{\nmref}[2]{\raisebox{1.1em}[0pt][0pt]{\hypertarget{ref:#1}{}}#2\par}
\newcommand{\nmdoi}[1]{\href{https://doi.org/#1}{\color{pblue}doi:#1}}
% Where the package is: its name and a link (its GitHub repository, once published)
\newcommand{\package}[2]{\href{#2}{\fn{#1.jl}}}

% Key facts at the top: source, missing in Julia, package, result
\newcommand{\factsheet}[4]{%
  \begin{tcolorbox}[enhanced, colback=pale, colframe=pale, boxrule=0pt, sharp corners,
    left=7pt, right=7pt, top=5pt, bottom=5pt, fontupper=\small]
  \begin{tabularx}{\linewidth}{@{}>{\leavevmode\bfseries\color{navy}}p{27mm}>{\leavevmode}X@{}}
  Tested against & #1\\
  Missing in Julia & #2\\
  Package & #3\\
  Result & #4
  \end{tabularx}
  \end{tcolorbox}}

% Boxes
\newtcolorbox{pbox}[2][]{enhanced, colback=pale, colframe=amber, boxrule=0pt,
  leftrule=3pt, sharp corners, left=7pt, right=7pt, top=6pt, bottom=6pt,
  fontupper=\small\raggedright, before upper={{\color{navy}\normalsize\bfseries #2}\par\smallskip}, #1}
\newtcolorbox{nbox}[2][]{enhanced, colback=navy, colframe=navy, boxrule=0pt,
  sharp corners, left=8pt, right=8pt, top=7pt, bottom=7pt, coltext=white,
  fontupper=\small\raggedright, before upper={{\color{amber}\normalsize\bfseries #2}\par\smallskip}, #1}

% Value tables: published beside Julia
\newenvironment{valuetable}[1]{%
  \par\small\tabularx{\linewidth}{#1}\toprule}{\bottomrule\endtabularx\par}
\newcommand{\ok}{{\color{agree}\faCheck}}
\newcommand{\no}{{\color{differ}\faTimes}}

% End of the document: credit and licence
\newcommand{\nmcredit}{%
  \par\vfill
  \begin{minipage}[b]{0.78\linewidth}\footnotesize\color{soft}
  A Pollis method, tested against the published tables on the published data.\par
  \textcopyright\ \vignetteauthor\enspace·\enspace Licensed under \vignettelicense\enspace·\enspace\vignettedate
  \end{minipage}\hfill
  \begin{minipage}[b]{0.18\linewidth}\raggedleft\pollislogo{9mm}\end{minipage}}
.
\documentclass[10pt,a4paper]{article}

\usepackage[a4paper,left=18mm,right=18mm,top=15mm,bottom=17mm,footskip=8mm]{geometry}
\usepackage{iftex}
\ifPDFTeX\usepackage[T1]{fontenc}\fi
\usepackage[default,semibold]{sourcesanspro}
\usepackage{textcomp}
\usepackage[italic]{mathastext}
\usepackage{amsmath}
\usepackage{microtype}
\usepackage[table]{xcolor}
\usepackage{tikz}
\usepackage[most]{tcolorbox}
\usepackage{fontawesome5}
\usepackage{tabularx,booktabs}
\usepackage{enumitem}
\usepackage{fancyhdr}
\usepackage[hidelinks]{hyperref}

% ---------------------------------------------------------------- colours (as in the brochure)
\definecolor{navy}{HTML}{1C3A5E}
\definecolor{pblue}{HTML}{3B78B5}
\definecolor{amber}{HTML}{F2B33D}
\definecolor{deepamber}{HTML}{D9902A}
\definecolor{pale}{HTML}{EEF3F9}
\definecolor{ink}{HTML}{2A2F36}
\definecolor{logonavy}{HTML}{173D63}
\definecolor{logodark}{HTML}{C98A1E}
\definecolor{logoamber}{HTML}{F2B446}
\definecolor{logocream}{HTML}{FEEFD2}
\definecolor{soft}{HTML}{5E6773}
\definecolor{agree}{HTML}{2F7D4F}
\definecolor{differ}{HTML}{B23B3B}
\color{ink}

\setlength{\parindent}{0pt}
\setlength{\parskip}{4pt}
\setlist{nosep,leftmargin=*,itemsep=2pt}
\renewcommand{\arraystretch}{1.15}

\pagestyle{fancy}
\fancyhf{}
\renewcommand{\headrulewidth}{0pt}
\fancyfoot[L]{\footnotesize\color{soft}\textbf{Pollis}\enspace Article Methods\enspace·\enspace\nmtheme}
\fancyfoot[R]{\footnotesize\color{soft}\thepage}
\newcommand{\nmtheme}{}
\newcommand{\vignetteauthor}{}
\newcommand{\vignettedate}{\today}
\newcommand{\vignettelicense}{\href{https://creativecommons.org/licenses/by-sa/4.0/}{CC BY-SA 4.0}}

% The Pollis logo, redrawn so it prints sharply at any size
\newcommand{\pollislogo}[1]{%
  \begin{tikzpicture}[x=#1/512,y=-#1/512,baseline=(current bounding box.center)]
    \fill[logonavy] (0,0) rectangle (512,512);
    \fill[logodark] (36,422) -- (104,256) -- (172,422) -- cycle;
    \fill[logoamber] (104,422) -- (203,180) -- (302,422) -- cycle;
    \fill[logocream] (203,422) -- (339,89) -- (476,422) -- cycle;
  \end{tikzpicture}}

% Top of the first page: logo, name; then the kicker, title and a one-line promise
\newcommand{\nmhead}[3]{%
  \renewcommand{\nmtheme}{#1}%
  \thispagestyle{fancy}%
  \begin{minipage}[c]{13mm}\pollislogo{11mm}\end{minipage}%
  \begin{minipage}[c]{\dimexpr\linewidth-13mm\relax}
    {\color{navy}\Large\bfseries Pollis}\enspace{\color{soft}Article Methods}\hfill
    {\color{deepamber}\small\bfseries\MakeUppercase{#1}}
  \end{minipage}\par\vspace{7mm}
  {\color{navy}\fontsize{22}{26}\selectfont\bfseries #2}\par\vspace{4pt}
  {\color{soft}\large #3}\par\vspace{5mm}}

\newcommand{\nmsection}[1]{\par\vspace{5pt}{\color{navy}\large\bfseries #1}\par\vspace{1pt}}
\newcommand{\nmsub}[1]{\par\vspace{2pt}{\color{deepamber}\bfseries #1}\par}
\newcommand{\gt}{\,{\color{deepamber}\bfseries ›}\,}
\newcommand{\menupath}[1]{{\bfseries\color{navy}#1}}
\newcommand{\fn}[1]{\texttt{\small #1}}
% Citations: \nmcite{key}{Breiman (2001)} links to the entry \nmref{key}{...} in the References
\newcommand{\nmcite}[2]{\hyperlink{ref:#1}{\color{pblue}#2}}
\newcommand{\nmref}[2]{\raisebox{1.1em}[0pt][0pt]{\hypertarget{ref:#1}{}}#2\par}
\newcommand{\nmdoi}[1]{\href{https://doi.org/#1}{\color{pblue}doi:#1}}
% Where the package is: its name and a link (its GitHub repository, once published)
\newcommand{\package}[2]{\href{#2}{\fn{#1.jl}}}

% Key facts at the top: source, missing in Julia, package, result
\newcommand{\factsheet}[4]{%
  \begin{tcolorbox}[enhanced, colback=pale, colframe=pale, boxrule=0pt, sharp corners,
    left=7pt, right=7pt, top=5pt, bottom=5pt, fontupper=\small]
  \begin{tabularx}{\linewidth}{@{}>{\leavevmode\bfseries\color{navy}}p{27mm}>{\leavevmode}X@{}}
  Tested against & #1\\
  Missing in Julia & #2\\
  Package & #3\\
  Result & #4
  \end{tabularx}
  \end{tcolorbox}}

% Boxes
\newtcolorbox{pbox}[2][]{enhanced, colback=pale, colframe=amber, boxrule=0pt,
  leftrule=3pt, sharp corners, left=7pt, right=7pt, top=6pt, bottom=6pt,
  fontupper=\small\raggedright, before upper={{\color{navy}\normalsize\bfseries #2}\par\smallskip}, #1}
\newtcolorbox{nbox}[2][]{enhanced, colback=navy, colframe=navy, boxrule=0pt,
  sharp corners, left=8pt, right=8pt, top=7pt, bottom=7pt, coltext=white,
  fontupper=\small\raggedright, before upper={{\color{amber}\normalsize\bfseries #2}\par\smallskip}, #1}

% Value tables: published beside Julia
\newenvironment{valuetable}[1]{%
  \par\small\tabularx{\linewidth}{#1}\toprule}{\bottomrule\endtabularx\par}
\newcommand{\ok}{{\color{agree}\faCheck}}
\newcommand{\no}{{\color{differ}\faTimes}}

% End of the document: credit and licence
\newcommand{\nmcredit}{%
  \par\vfill
  \begin{minipage}[b]{0.78\linewidth}\footnotesize\color{soft}
  A Pollis method, tested against the published tables on the published data.\par
  \textcopyright\ \vignetteauthor\enspace·\enspace Licensed under \vignettelicense\enspace·\enspace\vignettedate
  \end{minipage}\hfill
  \begin{minipage}[b]{0.18\linewidth}\raggedleft\pollislogo{9mm}\end{minipage}}
.
\documentclass[10pt,a4paper]{article}

\usepackage[a4paper,left=18mm,right=18mm,top=15mm,bottom=17mm,footskip=8mm]{geometry}
\usepackage{iftex}
\ifPDFTeX\usepackage[T1]{fontenc}\fi
\usepackage[default,semibold]{sourcesanspro}
\usepackage{textcomp}
\usepackage[italic]{mathastext}
\usepackage{amsmath}
\usepackage{microtype}
\usepackage[table]{xcolor}
\usepackage{tikz}
\usepackage[most]{tcolorbox}
\usepackage{fontawesome5}
\usepackage{tabularx,booktabs}
\usepackage{enumitem}
\usepackage{fancyhdr}
\usepackage[hidelinks]{hyperref}

% ---------------------------------------------------------------- colours (as in the brochure)
\definecolor{navy}{HTML}{1C3A5E}
\definecolor{pblue}{HTML}{3B78B5}
\definecolor{amber}{HTML}{F2B33D}
\definecolor{deepamber}{HTML}{D9902A}
\definecolor{pale}{HTML}{EEF3F9}
\definecolor{ink}{HTML}{2A2F36}
\definecolor{logonavy}{HTML}{173D63}
\definecolor{logodark}{HTML}{C98A1E}
\definecolor{logoamber}{HTML}{F2B446}
\definecolor{logocream}{HTML}{FEEFD2}
\definecolor{soft}{HTML}{5E6773}
\definecolor{agree}{HTML}{2F7D4F}
\definecolor{differ}{HTML}{B23B3B}
\color{ink}

\setlength{\parindent}{0pt}
\setlength{\parskip}{4pt}
\setlist{nosep,leftmargin=*,itemsep=2pt}
\renewcommand{\arraystretch}{1.15}

\pagestyle{fancy}
\fancyhf{}
\renewcommand{\headrulewidth}{0pt}
\fancyfoot[L]{\footnotesize\color{soft}\textbf{Pollis}\enspace Article Methods\enspace·\enspace\nmtheme}
\fancyfoot[R]{\footnotesize\color{soft}\thepage}
\newcommand{\nmtheme}{}
\newcommand{\vignetteauthor}{}
\newcommand{\vignettedate}{\today}
\newcommand{\vignettelicense}{\href{https://creativecommons.org/licenses/by-sa/4.0/}{CC BY-SA 4.0}}

% The Pollis logo, redrawn so it prints sharply at any size
\newcommand{\pollislogo}[1]{%
  \begin{tikzpicture}[x=#1/512,y=-#1/512,baseline=(current bounding box.center)]
    \fill[logonavy] (0,0) rectangle (512,512);
    \fill[logodark] (36,422) -- (104,256) -- (172,422) -- cycle;
    \fill[logoamber] (104,422) -- (203,180) -- (302,422) -- cycle;
    \fill[logocream] (203,422) -- (339,89) -- (476,422) -- cycle;
  \end{tikzpicture}}

% Top of the first page: logo, name; then the kicker, title and a one-line promise
\newcommand{\nmhead}[3]{%
  \renewcommand{\nmtheme}{#1}%
  \thispagestyle{fancy}%
  \begin{minipage}[c]{13mm}\pollislogo{11mm}\end{minipage}%
  \begin{minipage}[c]{\dimexpr\linewidth-13mm\relax}
    {\color{navy}\Large\bfseries Pollis}\enspace{\color{soft}Article Methods}\hfill
    {\color{deepamber}\small\bfseries\MakeUppercase{#1}}
  \end{minipage}\par\vspace{7mm}
  {\color{navy}\fontsize{22}{26}\selectfont\bfseries #2}\par\vspace{4pt}
  {\color{soft}\large #3}\par\vspace{5mm}}

\newcommand{\nmsection}[1]{\par\vspace{5pt}{\color{navy}\large\bfseries #1}\par\vspace{1pt}}
\newcommand{\nmsub}[1]{\par\vspace{2pt}{\color{deepamber}\bfseries #1}\par}
\newcommand{\gt}{\,{\color{deepamber}\bfseries ›}\,}
\newcommand{\menupath}[1]{{\bfseries\color{navy}#1}}
\newcommand{\fn}[1]{\texttt{\small #1}}
% Citations: \nmcite{key}{Breiman (2001)} links to the entry \nmref{key}{...} in the References
\newcommand{\nmcite}[2]{\hyperlink{ref:#1}{\color{pblue}#2}}
\newcommand{\nmref}[2]{\raisebox{1.1em}[0pt][0pt]{\hypertarget{ref:#1}{}}#2\par}
\newcommand{\nmdoi}[1]{\href{https://doi.org/#1}{\color{pblue}doi:#1}}
% Where the package is: its name and a link (its GitHub repository, once published)
\newcommand{\package}[2]{\href{#2}{\fn{#1.jl}}}

% Key facts at the top: source, missing in Julia, package, result
\newcommand{\factsheet}[4]{%
  \begin{tcolorbox}[enhanced, colback=pale, colframe=pale, boxrule=0pt, sharp corners,
    left=7pt, right=7pt, top=5pt, bottom=5pt, fontupper=\small]
  \begin{tabularx}{\linewidth}{@{}>{\leavevmode\bfseries\color{navy}}p{27mm}>{\leavevmode}X@{}}
  Tested against & #1\\
  Missing in Julia & #2\\
  Package & #3\\
  Result & #4
  \end{tabularx}
  \end{tcolorbox}}

% Boxes
\newtcolorbox{pbox}[2][]{enhanced, colback=pale, colframe=amber, boxrule=0pt,
  leftrule=3pt, sharp corners, left=7pt, right=7pt, top=6pt, bottom=6pt,
  fontupper=\small\raggedright, before upper={{\color{navy}\normalsize\bfseries #2}\par\smallskip}, #1}
\newtcolorbox{nbox}[2][]{enhanced, colback=navy, colframe=navy, boxrule=0pt,
  sharp corners, left=8pt, right=8pt, top=7pt, bottom=7pt, coltext=white,
  fontupper=\small\raggedright, before upper={{\color{amber}\normalsize\bfseries #2}\par\smallskip}, #1}

% Value tables: published beside Julia
\newenvironment{valuetable}[1]{%
  \par\small\tabularx{\linewidth}{#1}\toprule}{\bottomrule\endtabularx\par}
\newcommand{\ok}{{\color{agree}\faCheck}}
\newcommand{\no}{{\color{differ}\faTimes}}

% End of the document: credit and licence
\newcommand{\nmcredit}{%
  \par\vfill
  \begin{minipage}[b]{0.78\linewidth}\footnotesize\color{soft}
  A Pollis method, tested against the published tables on the published data.\par
  \textcopyright\ \vignetteauthor\enspace·\enspace Licensed under \vignettelicense\enspace·\enspace\vignettedate
  \end{minipage}\hfill
  \begin{minipage}[b]{0.18\linewidth}\raggedleft\pollislogo{9mm}\end{minipage}}

\renewcommand{\vignetteauthor}{Antonio Saragga Seabra}
\renewcommand{\vignettedate}{11 November 2026}
\renewcommand{\vignettelicense}{the MIT License}

\begin{document}

\nmhead{High-frequency volatility}
  {Two-Scales Realized Volatility}
  {The daily variance of a price observed every second through market noise: the two-scales estimator of Zhang,
   Mykland and Aït-Sahalia (2005), with their Monte Carlo study reproduced.}

\factsheet
  {\nmcite{zma2005}{Zhang, Mykland and Aït-Sahalia (2005)}, \emph{Journal of the American Statistical Association}:
   Table~2 (p.\,1405), the five estimators of the integrated variance on 25,000 simulated days of the Heston model;
   no data. Also a known volatility, and \fn{rTSCov} of the R package highfrequency 1.0.3.}
  {the two-scales estimator as the article writes it. \nmcite{bk2022}{Baumann and Klymak (2022)}, in
   HighFrequencyCovariance.jl, apply its correction to volatilities, the square roots of the realized variances, not to
   the variances.}
  {\fn{TwoScalesVolatility.jl}}
  {43 of the 45 printed values tested. The 4 that the theory gives without simulation: 3 to every printed digit, 1 off
   by one in the last. The 39 simulated: all within the noise of a second run, the largest gap 2.49 standard errors.}

\nmsection{The problem}
The sum of a day's squared returns estimates its integrated variance $\langle X,X\rangle_T=\int_0^T\sigma_t^2\,dt$
when prices carry no error. With noise, $Y_{t_i}=X_{t_i}+\varepsilon_{t_i}$, it grows like $2n\,E\varepsilon^2$ with
the $n$ returns: every second, it measures the noise. Five-minute returns throw most of the data away and keep a bias;
the article combines two time scales to remove it with all the data.

\nmsection{The method}
Table~2 ranks five estimators: every tick, $[Y,Y]^{(\mathrm{all})}$ (fifth-best); every five minutes (fourth-best);
sparse returns at the optimal step (third-best); the average $[Y,Y]^{(\mathrm{avg})}$ of the realized variances of $K$
regular sub-grids of the optimal size, equation~(32) (second-best); and the two-scales estimator (first-best),
equations~(55) and~(64):
\[
\widehat{\langle X,X\rangle}_T=\Bigl(1-\frac{\bar n}{n}\Bigr)^{-1}\Bigl([Y,Y]^{(\mathrm{avg})}_T-\frac{\bar n}{n}\,[Y,Y]^{(\mathrm{all})}_T\Bigr),
\qquad \bar n=\frac{n-K+1}{K},\qquad K=c^*n^{2/3},\quad
c^*=\Bigl(\frac{T}{12\,(E\varepsilon^2)^2}\int_0^T\sigma_t^4\,dt\Bigr)^{-1/3}.
\]
The noise variance is $[Y,Y]^{(\mathrm{all})}/(2n)$, equation~(22), or $([Y,Y]^{(\mathrm{all})}-[Y,Y]^{(\mathrm{avg})})/(2(n-\bar n))$, equation~(65).

\fn{realized\_variance}, \fn{sparse\_variance}, \fn{average\_variance} and \fn{two\_scales\_variance} return the four
statistics from a day's log prices, \fn{noise\_variance} equation~(22) or, given $K$, (65), and \fn{optimal\_sparse},
\fn{optimal\_average} and \fn{optimal\_c} the tuning. \fn{HestonDesign} holds the design: the model of
\nmcite{heston1993}{Heston (1993)}, $\mu=0.05$, $\kappa=5$, $\alpha=0.04$, $\gamma=0.5$, $\rho=-0.5$, noise of standard
deviation 0.0005, 23,400 one-second prices in a day of $T=1/252$ years. \fn{simulate\_day} draws a day by the Euler
scheme, \fn{table2} returns Table~2 as a $9\times5$ matrix and \fn{monte\_carlo\_table2} adds its standard errors;
\fn{simulate\_constant\_day} and \fn{recovery} test a known volatility. Four choices the article leaves unsaid are
inferred: the variance starts from its stationary Gamma law (shape 1.6, scale 0.025), which the five-minute asymptotic
variance fits; each day is tuned with its own quarticity and the true noise variance; the sparse and averaged
estimators cover the whole day, as the small-sample biases show; variances are taken within $50\times5$ bins of the
days by integrated variance, then quarticity.

\nmsection{What the tests show}
A known volatility of 1\% a day (integrated variance $10^{-4}$), noise of 0 to 20 basis points, $K=c^*n^{2/3}$ (2, 35,
87, 219); means over 400 days. Its sub-grids leave out $K-1$ ticks on average, so the estimator is
expected to give $10^{-4}(1-K/(n+1))$.
{\footnotesize
\begin{valuetable}{@{}lXrrrrc@{}}
\textbf{Value} & & \textbf{0 bp} & \textbf{5 bp} & \textbf{10 bp} & \textbf{20 bp} & \\\midrule
Integrated variance & Expected & 9.999e-5 & 9.985e-5 & 9.963e-5 & 9.906e-5 & \\
 & Julia & 9.989e-5 & 9.947e-5 & 9.961e-5 & 1.0036e-4 & \ok\\
Noise variance, (65) & Known & 0 & 2.500e-7 & 1.000e-6 & 4.000e-6 & \\
 & Julia & 2.2e-12 & 2.503e-7 & 1.001e-6 & 3.998e-6 & \ok\\
\end{valuetable}}
Against \fn{rTSCov} of R highfrequency 1.0.3, one day each at 5 bp ($K=35$), 10 bp ($K=87$), 20 bp ($K=219$) and 5 bp
($K=300$), the relative differences are $2.1\times10^{-7}$, $9.1\times10^{-7}$, $3.4\times10^{-6}$ and
$1.8\times10^{-7}$, under the $10^{-5}$ tested: highfrequency counts the $n+1$ prices where the article counts the $n$
returns, in $\bar n$ and $\bar n/n$.

Table~2, p.\,1405, published above Julia (25,000 days). A simulated value is compared as Pollis compares random values:
the gap beyond the rounding of the printed value, in standard errors of the difference of two runs, $z$, must be
under~3. The cells in bold, given by the theory, are tested to the digits printed; the two in grey are not tested.
{\footnotesize
\begin{valuetable}{@{}Xlrrrrrrc@{}}
\textbf{Row} & & \textbf{Fifth} & \textbf{Fourth} & \textbf{Third} & \textbf{Second} & \textbf{First} & \textbf{Largest $z$} & \\\midrule
Small-sample bias & Published & 1.1699e-2 & 3.89e-5 & 2.18e-5 & 1.926e-5 & 2e-8 & & \\
 & Julia & 1.1702e-2 & 3.91e-5 & 2.19e-5 & 1.930e-5 & -7e-8 & 1.85 & \ok\\
Asymptotic bias & Published & \textbf{1.1700e-2} & \textbf{3.90e-5} & 2.20e-5 & 1.927e-5 & \textbf{0} & & \\
 & Julia & \textbf{1.1700e-2} & \textbf{3.90e-5} & 2.19e-5 & 1.912e-5 & \textbf{0} & 1.63 & \ok\\
Small-sample variance & Published & 1.791e-8 & 1.4414e-9 & 1.59e-9 & 9.41e-10 & 9e-11 & & \\
 & Julia & 1.783e-8 & 1.4099e-9 & 1.62e-9 & 9.62e-10 & 9e-11 & 1.18 & \ok\\
Asymptotic variance & Published & 1.788e-8 & 1.4409e-9 & 1.58e-9 & 9.37e-10 & 8e-11 & & \\
 & Julia & 1.787e-8 & 1.4280e-9 & 1.58e-9 & 9.35e-10 & 8e-11 & 1.47 & \ok\\
Small-sample RMSE & Published & 1.1699e-2 & 5.437e-5 & 4.543e-5 & 3.622e-5 & 9.4e-6 & & \\
 & Julia & 1.1702e-2 & 5.423e-5 & 4.580e-5 & 3.653e-5 & 9.4e-6 & 2.49 & \ok\\
Asymptotic RMSE & Published & \textbf{1.1700e-2} & 5.442e-5 & 4.546e-5 & 3.618e-5 & 8.9e-6 & & \\
 & Julia & \textbf{1.1701e-2} & 5.430e-5 & 4.532e-5 & 3.606e-5 & 8.9e-6 & 0.85 & \ok\\
Relative bias & Published & 182 & .61 & .18 & .15 & -.00045 & & \\
 & Julia & 181 & .60 & .17 & .15 & -.00147 & 1.76 & \ok\\
Relative variance & Published & {\color{soft}82502} & 1.15 & .11 & .053 & .0043 & & \\
 & Julia & {\color{soft}123200} & 1.06 & .11 & .056 & .0042 & 1.48 & \ok\\
Relative RMSE & Published & {\color{soft}340} & 1.24 & .37 & .28 & .065 & & \\
 & Julia & {\color{soft}395} & 1.19 & .37 & .28 & .065 & 0.13 & \ok\\
\end{valuetable}}

\nmsub{Where they differ}
\begin{itemize}
  \item \textbf{Asymptotic RMSE, fifth-best:} $\sqrt{0.0117^2+1.787\times10^{-8}}=0.0117008$ rounds to 1.1701e-2, one
    unit above the printed 1.1700e-2.
  \item \textbf{Simulated values} differ in their digits, as two runs must: the five-minute small-sample variance by
    315 units of its last digit, at $z=0.98$. The first-best biases carry the $-K/(n+1)$ of the ticks left out.
  \item \textbf{Not tested:} the fifth-best relative variance and RMSE. The integrated variance has a Gamma law of shape
    1.6, so the mean of its inverse square is infinite and they never settle (standard errors 79,700 and 101).
  \item \textbf{Known volatility:} every mean within 3 standard errors of what is expected (the largest gap 2.34, the
    noise variance at 5 bp), the integrated variance within 0.6\% of $10^{-4}$.
\end{itemize}

\nmsection{Still open}
Figures~1 and~2 are not reproduced. The article does not say whether its optimal estimators are tuned with true or
estimated values; the true ones are used here.

\nmsection{Further considerations}
\begin{itemize}
  \item \textbf{A limit.} The noise is independent of the price and from tick to tick, the prices equally spaced; a
    jump in the day enters the estimate.
  \item \textbf{An extension.} More than two scales, or dependent noise, each tested against its own article; the
    correction applied to variances could be offered to HighFrequencyCovariance.jl.
\end{itemize}

\nmsection{References}
{\footnotesize
\nmref{zma2005}{Zhang, L., Mykland, P. A. and Aït-Sahalia, Y. (2005). A tale of two time scales: determining integrated
  volatility with noisy high-frequency data. \emph{Journal of the American Statistical Association} 100(472),
  1394--1411. \nmdoi{10.1198/016214505000000169}}
\nmref{bk2022}{Baumann, S. and Klymak, M. (2022). HighFrequencyCovariance: a Julia package for estimating covariance
  matrices using high frequency financial data. \emph{Journal of Statistical Software} 103(14).
  \nmdoi{10.18637/jss.v103.i14}}
\nmref{heston1993}{Heston, S. L. (1993). A closed-form solution for options with stochastic volatility with
  applications to bond and currency options. \emph{Review of Financial Studies} 6(2), 327--343.
  \nmdoi{10.1093/rfs/6.2.327}}
}

\nmcredit
\end{document}
